%! Bias in Samples and Surveys — Unit 1, Lesson 1.4 — Guided Notes & Practice
% PILOT SHAPE (2026-09-06): modeled on AP Statistics lesson 1.4 minus its AP Practice page.
% THE in-class document. The notes carry the whole gradual release, closed by a SPOKEN debrief.
% They run SHORT on purpose: the period ends with students starting the homework, which is
% where the individual practice lives.
%   I DO   : the vocabulary box and notes sections 1-2.                              (~20 min)
%   WE DO  : the practicebox ("Guided Practice") — Cedar Ridge Apartments, together.  (~15 min)
%   YOU DO : nothing on this page — ../homework/ IS the individual practice, started in the
%            last ten minutes of the period.
% PAGE PLAN: p1 = vocab box + hook (no objective box — the cover carries the targets);
% p2 = section 1; p3 = section 2; p4 = the Guided Practice alone. Every fill-in cluster opens
% with a \wordbank{}. No closing practice set, no extension.
% Vocabulary is GIVEN here, not discovered. THE CRUX is the second half of section 2: "the
% sample was too small" — a larger sample makes an estimate more PRECISE, not more CORRECT;
% size shrinks variability and does nothing to direction (PS.DC.2c, AFDA.DA.2e). The second
% half of section 1 carries the second trap: undercoverage and nonresponse fused — the test is
% COULD THAT PERSON HAVE BEEN CHOSEN?
%
% THE DATA
%   Harbor Point Community Pool (warm-up + notes): 1,500 members; census truth 300/1500 = 20%
%     mainly swim laps. Tuesday-evening sample 27/75 = 36% -> 540 predicted, off by 240.
%     Three honest samples 19/22/19 average 20%; three Tuesday samples 35/36/37 average 36%.
%     Ten Tuesdays: 270/750 = 36% — the same 16-point miss. Question order: 62% vs 41% = 21.
%   Guided Practice, Cedar Ridge Apartments (the old activity's context): 600 households ALL
%     mailed a card; 150 return (25%); 96/150 = 64% -> 384 predicted; census 240/600 = 40%,
%     so 24 points and 144 households too many; re-run 192/300 = 64% again.
%
% PAGE LOCKSTEP with notes_key/: the key is GENERATED from this file — every \blank{} becomes
% an \ans{}, every \par\writelines{n} becomes n \ansline{} calls, every \vterm{} a
% \vtermans{}{}. Every \begin{work} block is byte-identical. Blank widths are sized to their
% answers so the two files wrap alike.
\documentclass[12pt]{article}
\usepackage{saar-article}
\usepackage{saar-boxes}
\newcommand{\TallMath}[1]{$\displaystyle #1\rule[-1.4em]{0pt}{3.2em}$}
% Word bank strip for every fill-in cluster (user request 2026-09-06). Defined per document;
% the key inherits it and prints the same strip, so heights match.
\newcommand{\wordbank}[1]{\par\vspace{2pt}\noindent\fcolorbox{goldacc}{goldbg}{\parbox{\dimexpr\linewidth-2\fboxsep-2\fboxrule\relax}{\small\textbf{Word bank:} #1}}\par\vspace{4pt}}
% Vocabulary rows: a FIXED-HEIGHT row carrying the term label and OPEN writing space —
% no underline beside the term and no rule line (user correction 2026-09-06: the black
% answer line crowded the row; students write beside and below the term). The key fills
% exactly the same height, so the box cannot drift blank vs. keyed. saar-article's
% \termblank adds an inline blank and a rule line, so the pair is defined here (the key is
% generated from this file and inherits both). 2.3cm per row gives students three
% handwritten lines of room; keep key definitions to two lines.
\newcommand{\termrowheight}{2.3cm}
\newlength{\termrowinset}\setlength{\termrowinset}{7pt}
\newcommand{\vterm}[1]{%
  \par\noindent\begin{minipage}[t][\termrowheight][t]{\linewidth}%
    \vspace*{\termrowinset}%
    \textbf{\textcolor{cerulean}{#1:}}%
  \end{minipage}\par}
\newcommand{\vtermans}[2]{%
  \par\noindent\begin{minipage}[t][\termrowheight][t]{\linewidth}%
    \vspace*{\termrowinset}%
    \textbf{\textcolor{cerulean}{#1:}}\quad\ans{#2}%
  \end{minipage}\par}

\begin{document}

\pageheader{Unit 1, Lesson 1.4}{Guided Notes \& Practice}
% NAMESTRIP: no \namedateperiod here. The student writes their name once, on the
% cover this component is stapled behind.

\vspace{0.15in}

\begin{vocabbox}
\small
Fill in each term \textbf{as we name it} in the notes below.
\vterm{Bias}
\vterm{Sampling bias}
\vterm{Undercoverage}
\vterm{Nonresponse}
\vterm{Response bias}
\end{vocabbox}

\boxguard[12]
\begin{hookbox}
\small
Last class the board threw out the Tuesday-evening study: $27$ of the $75$ members asked said
they mainly swim laps --- $36\%$ --- while the office's own records say $20\%$. The staff
member says the problem was \emph{size}. She will run the same Tuesday-evening survey for
ten weeks and ask $750$ members instead of $75$.

\vspace{3pt}
\textbf{Will her new number land closer to the truth, $20\%$?} Circle one: \quad yes \qquad
no \qquad it depends \quad --- and say why you think so. We vote again at the end of
section~2.
\par\writelines{2}
\end{hookbox}

% ── I DO: section 1 ──────────────────────────────────────────────────────────
% Page 2: section 1 (page 1 is the vocab box + hook).
\newpage
\begin{notesbox}{1. Bias Is a Direction, Not Bad Luck}
\small
Lesson 1.2 said that honest random samples disagree with each other --- that is
\textbf{sampling variability}, and nobody is wrong. Those samples scatter \emph{around} the
truth. Today's studies are wrong the \emph{same way every time}.
\wordbank{direction \;\textbullet\; 20\% \;\textbullet\; 36\% \;\textbullet\; Tuesday \;\textbullet\; too high \;\textbullet\; too low \;\textbullet\; lap swimmers \;\textbullet\; undercoverage \;\textbullet\; nonresponse \;\textbullet\; could \;\textbullet\; could not}
A study is \textbf{biased} when its procedure \emph{systematically} overestimates or
underestimates the population value. Every sample it takes leans the same \blank{1.9cm}.
The office's records say $20\%$ of all $1{,}500$ members mainly swim laps (the warm-up).
Here are two ways of taking three samples of $75$:
\nopagebreak

\begin{center}
\renewcommand{\arraystretch}{1.3}
\begin{tabularx}{\linewidth}{@{} X c c c p{1.9cm} @{}}
\toprule
\textbf{How the $75$ were chosen} & \textbf{1st} & \textbf{2nd} & \textbf{3rd} & \textbf{Average} \\
\midrule
At random from the full membership list & $19\%$ & $22\%$ & $19\%$ & \blank{1.5cm} \\
From whoever arrived Tuesday evening    & $35\%$ & $36\%$ & $37\%$ & \blank{1.5cm} \\
\bottomrule
\end{tabularx}
\end{center}

\vspace{2pt}
Both rows disagree with themselves --- that is variability, and it is unavoidable. Only the
\blank{1.9cm} row misses the \emph{truth} every single time, and it misses \blank{1.9cm}.
Tuesday evening is lap-swim time, so \blank{2.8cm} are over-counted. \textit{Consistent
is not the same as correct: the tighter row is the wrong one.}

\vspace{5pt}
\textbf{\textcolor{cerulean}{Sampling Bias --- Who Never Had a Chance.}}
\textbf{Sampling bias} gets in when some members of the population are systematically more
likely to be chosen than others. It arrives two ways, both before anyone speaks.
\textbf{Undercoverage:} part of the population is left off the sampling frame, so those
members can never be chosen. \textbf{Nonresponse:} members are chosen properly but never
answer, so the data is missing exactly the people who stayed silent.
\nopagebreak

\begin{center}
\renewcommand{\arraystretch}{1.25}
\begin{tabularx}{\linewidth}{@{} X p{3.9cm} p{2.9cm} @{}}
\toprule
\textbf{The plan} & \textbf{Who is missing?} & \textbf{Which bias?} \\
\midrule
Ask the members who arrive Tuesday evening. & morning and weekend swimmers & \blank{2.6cm} \\
Email the survey only to members with an address on file. & members with no email on file & \blank{2.6cm} \\
Post the survey on the pool's social media page. & members who do not follow the page & \blank{2.6cm} \\
Mail a card to all $1{,}500$ members; $200$ mail it back. & the $1{,}300$ who never mailed it back & \blank{2.6cm} \\
\bottomrule
\end{tabularx}
\end{center}

\vspace{3pt}
\begin{center}
\fcolorbox{redacc}{redbg}{\parbox{0.94\linewidth}{\small
\textbf{Careful --- the last row's frame is perfect.} All $1{,}500$ got a card, and the
study is still broken. The test that separates the two: \emph{could that person have been
chosen?} A morning swimmer \blank{1.8cm} have been chosen on Tuesday evening ---
undercoverage. A silent card holder \blank{1.4cm} have been --- nonresponse. Both leave
people out, and neither one shows up in your arithmetic.}}
\end{center}

\end{notesbox}

% Page 3: section 2.
\newpage
\begin{notesbox}{2. Response Bias --- After the Right Person Is Chosen}
\small
\textbf{Response bias} is anything about the survey itself --- who is asking, how it is worded,
what order it comes in --- that pushes a respondent away from their true answer. It happens
\emph{after} the right person has been chosen. The standard names seven kinds; read them once,
then diagnose.
\wordbank{social desirability \;\textbullet\; neutral responding \;\textbullet\; acquiescence \;\textbullet\; demand \;\textbullet\; question order \;\textbullet\; 36\% \;\textbullet\; 16 \;\textbullet\; precise \;\textbullet\; correct \;\textbullet\; variability \;\textbullet\; direction \;\textbullet\; undercoverage \;\textbullet\; nonresponse}
\nopagebreak
\begin{center}
\renewcommand{\arraystretch}{1.1}
\begin{tabularx}{\linewidth}{@{} >{\bfseries}p{4.1cm} X @{}}
\toprule
\textbf{\normalfont Kind} & \textbf{What the respondent does} \\
\midrule
Demand bias & Guesses what the study wants to hear and answers to match. \\
Social desirability & Gives the answer that looks good to other people. \\
Dissent bias & Answers negatively to everything, often without understanding the questions. \\
Acquiescence bias & Agrees with everything --- every answer positive. \\
Extreme responses & Always ``strongly agree'' or ``strongly disagree.'' \\
Neutral responding & Always the middle, no-opinion choice. \\
Question order bias & An earlier question changes the answer to a later one. \\
\bottomrule
\end{tabularx}
\end{center}

\nopagebreak
\begin{center}
\renewcommand{\arraystretch}{1.15}
\begin{tabularx}{\linewidth}{@{} X p{3.4cm} @{}}
\toprule
\textbf{At the pool} & \textbf{Kind} \\
\midrule
The lifeguard hands it out and waits; last week's complainers rate the staff ``excellent.'' & \blank{3.1cm} \\
A member marks ``no opinion'' on all fifteen questions. & \blank{3.1cm} \\
A member agrees with every statement, even two that contradict each other. & \blank{3.1cm} \\
The survey is titled \emph{``Help us show the pool needs longer hours.''} & \blank{3.1cm} \\
Asked \emph{``Are the lap lanes overcrowded?''} first, $62\%$ say yes to a \$40,000 lane;
asked the money question first, $41\%$ do --- same words, $21$ points apart. & \blank{3.1cm} \\
\bottomrule
\end{tabularx}
\end{center}

\vspace{4pt}
\textbf{\textcolor{cerulean}{The Fix That Is Not a Fix.}}
Told her sample was too small, the staff member runs the same Tuesday-evening survey for ten
weeks. She asks $750$ members, and $270$ of them say they mainly swim laps. Show the work:

\begin{work}
  \dfrac{270}{750} &= 0.36 = 36\%
\end{work}

One Tuesday: $75$ asked, $36\%$, off by \blank{1.1cm} points. Ten Tuesdays: $750$ asked,
\blank{1.3cm}, off by \blank{1.1cm} points. Ten times the work, the same answer.
\textbf{Vote again.}

\vspace{3pt}
\begin{center}
\fcolorbox{cerulean}{frost}{\parbox{0.94\linewidth}{\small
\textbf{A larger sample makes an estimate more \blank{1.7cm}, not more \blank{1.7cm}.} Size
shrinks \blank{2.1cm}; nothing about size touches \blank{1.8cm}. The $750$ are still
Tuesday-evening lap swimmers --- a bigger biased sample is a more confident wrong answer. Only
a better \emph{procedure} removes bias, and each fix removes one: a complete frame removes
\blank{2.6cm}; calling back the silent removes \blank{2.4cm}; neutral wording and anonymous
answers remove response bias.}}
\end{center}

\end{notesbox}

% Page 4: the Guided Practice, alone.
\newpage
\begin{practicebox}
\small
\textbf{Cedar Ridge Apartments --- we work this one together.} Cedar Ridge has $600$
households. Management mails a parking survey to \emph{every} household; $150$ cards come
back, and $96$ of those say parking is a problem. Management is about to tell the owners that
Cedar Ridge needs a bigger lot.
\wordbank{undercoverage \;\textbullet\; nonresponse \;\textbullet\; yes \;\textbullet\; no \;\textbullet\; 24 \;\textbullet\; 144 \;\textbullet\; too high \;\textbullet\; too low \;\textbullet\; acquiescence \;\textbullet\; demand}
\textbf{(a)} What share of the households answered at all, and what share of \emph{those}
said parking is a problem?

\begin{work}
  \dfrac{150}{600} &= 0.25 = 25\% \\
  \dfrac{96}{150}  &= 0.64 = 64\%
\end{work}

\vspace{3pt}
\textbf{(b)} Every household got a card, so nobody was left off the frame --- this is not
\blank{2.8cm}. The $450$ households that threw the card away are \blank{2.6cm}. Could one
of them have been chosen? \blank{1cm}

\vspace{3pt}
\textbf{(c)} Management scales $64\%$ to all $600$ households. A year later the office asks
\emph{every} household, and $240$ of the $600$ say parking is a problem. Show both.

\begin{work}
  0.64 \times 600 &= 384 \text{ households} \\
  \dfrac{240}{600} &= 0.40 = 40\%
\end{work}

The mail-in study was off by \blank{1.2cm} percentage points --- \blank{1.4cm} households
--- and it was \blank{2cm}.

\vspace{3pt}
\textbf{(d)} Management says the first study just needed more cards. They mail again, get
$300$ back, and $192$ of those say parking is a problem: $192 \div 300 = 64\%$ again.
Two-bedroom households own the most cars \emph{and} mailed their cards back at twice the
rate of studio households. Explain why twice the cards did not help, and why the miss is in
the direction you named in (c).
\par\writelines{2}

\vspace{3pt}
\textbf{(e)} The card asked: \emph{``Don't you agree that Cedar Ridge's parking lot is
dangerously overcrowded?''} Name the response bias that wording invites: \blank{3cm}. Then
rewrite the question so it pushes neither way.
\par\writeline
\end{practicebox}

\end{document}
